\section{Discussion}

The experimental evaluations across the two distinct benchmarks expose why compactness cannot be judged by a single headline score. On the balanced four-class benchmark, multiple compact models reach approximately 98\% accuracy, and paired McNemar tests do not resolve a difference between LS-Net and ResNet-50 for the reported seeds. This supports a useful operating point at 0.91M parameters and 0.36 GFLOPs without implying that all compact models are equivalent or that added capacity has no benefit. On BraTS, the strongest model evaluated is instead VGG-16, and the gap between its 89.31\% slice F1 and LS-Net's 82.02\% is 7.29 points. Transformer-based architectures have since become competitive baselines for brain tumor analysis, but their multi-million-parameter designs make it worth measuring what compact CNNs can and cannot achieve at the edge. Different model rankings and class-level outcomes warn against using performance on a curated tumor-type benchmark as a proxy for HGG/LGG grading.

MCG-TTA changes the allocation of compute rather than the representational capacity of the learned classifier. It uses identity confidence to select a subset of slices and scale-view disagreement to decide whether all views are needed. The ablation makes clear that the value of the architecture-informed rule is conditional: its more favorable LGG preservation relative to some alternative gates is an observed result on 24 pooled LGG patient decisions. The residual miss rate, reduced patient F1 versus the unaugmented model, and absence of independent clinical validation rule out autonomous diagnostic claims.

Several limitations constrain the interpretation. BraTS grading is a separate training-and-testing problem, so these results measure robustness of model choice across tasks rather than invariance of one learned decision rule. Mask-guided selection of tumor-containing slices assumes information that may be unavailable when screening a new scan. Slice-level observations are correlated within subjects, and the minority-class patient count is small. The four-class table reports three seeds, but the detailed TTA ablation reports two; claims about the TTA gain should be limited to that ablation. McNemar's non-significant results are not equivalence proofs, and adaptive latency is modelled from forward-pass counts rather than directly measured on a clinical device. Finally, the present study does not measure lesion localization, segmentation, cloud privacy outcomes, calibration, or radiologist performance.

These boundaries point to concrete next evaluations: prospective subject-level grading, direct end-to-end latency on target hardware, uncertainty calibration, and studies in which slice selection does not rely on segmentation masks. They will determine whether the present efficiency and slice-level gains translate into a reliable clinical patient workflow.
